\documentclass[11pt]{article}
\usepackage{amsmath,amssymb,amsthm}
\usepackage{fullpage}
\usepackage{amsmath,amssymb,amsthm}
\usepackage[hidelinks]{hyperref}

\newcommand{\N}{\mathbb N}
\newcommand{\Z}{\mathbb Z}
\newcommand{\gp}{\mathrm{gp}}
\newcommand{\sat}{\mathrm{sat}}
\newcommand{\Hom}{\operatorname{Hom}}
\newcommand{\Sh}{\operatorname{Sh}}
\newcommand{\PSh}{\operatorname{PSh}}
\newcommand{\Ab}{\operatorname{Ab}}
\newcommand{\R}{\mathrm R}
\newcommand{\Jb}{J_{\mathrm b}}
\newcommand{\G}{\mathcal G}
\newcommand{\E}{\mathcal E}
\newcommand{\I}{\mathcal I}
\newcommand{\V}{\mathcal V}
\newcommand{\Lan}{\operatorname{Lan}}
\newcommand{\coker}{\operatorname{coker}}

\newtheorem{theorem}{Theorem}[section]
\newtheorem{proposition}[theorem]{Proposition}
\newtheorem{lemma}[theorem]{Lemma}
\newtheorem{corollary}[theorem]{Corollary}
\theoremstyle{definition}
\newtheorem{definition}[theorem]{Definition}
\theoremstyle{remark}
\newtheorem{remark}[theorem]{Remark}

\newcommand{\fptitle}{Review of Problem 1 -- ProofCouncil}
\newcommand{\reviewtag}{Reviewer 1}
\usepackage{fp-review}

\begin{document}
\maketitle

\section{Review}

This solution was clearly the best of those we reviewed.  It followed the same general lines as the human solution.  This solution made use of more abstract tools than the other solutions (it introduced an auxiliary topology), which made it conceptually simpler.  This solution also used a Mayer--Vietoris argument where the authors had constructed a contractible simplicial complex, which gave an additional conceptual simplification.  The writing is excellent overall, although there is room to tailor the presentation better for a human reader.  I did not identify any inadequate citations.

\paragraph{Correctness.}
\begin{itemize}
	\item  1: essentially flawless mathematical argument
	\item  0.5: minor errors that are easily corrected
	\item  0: blank solution, or acknowledged partial result
	\item  $-0.5$: major errors
	\item  $-1$: proposed proof makes no progress towards the result.
\end{itemize}
Note that points should be docked for wrong statements/ solutions.

\noindent Score: 1

\paragraph{Exposition.}
\begin{itemize}
	\item  1: well-written and clear solution.
	\item  0: argument that is possible to follow with serious effort
	\item  $-1$: very poor exposition
\end{itemize}

\noindent Score: 1

\paragraph{Attributions.}
\begin{itemize}
	\item  1: solutions contains fully appropriate attributions
		to related work and theorems being used.
	\item  0: moderate issues with missing attributions or hallucinated references
	\item  $-1$: blatant plagiarism or hallucinations
\end{itemize}

\noindent Score: 1

\newpage
\section{Submission with inline comments}


\begin{abstract}
Let \(P\) be an integral, saturated, sharp commutative monoid, let \(A_P\) be
the category of monoids under \(P\), equipped with the birhizomic topology on
\(A_P^{\mathrm{op}}\), and let \(B_P\subset A_P\) be the full subcategory of
the sharpened localizations \(P\to P[-x]^\sharp\), \(x\in P\), with the
trivial topology (all definitions are recalled in \S1).  The inclusion
induces an inverse image functor \(\pi^{-1}\) from presheaves on
\(B_P^{\mathrm{op}}\) to birhizomic sheaves, with right adjoint \(\pi_*\),
and the question answered here asks whether \(F\to\R\pi_*\pi^{-1}F\) is a
quasi-isomorphism for every presheaf of abelian groups \(F\) on
\(B_P^{\mathrm{op}}\).  We show that it is, for every \(P\), with no
finiteness hypothesis.  The proof replaces the birhizomic topology by a
finer auxiliary topology generated by binary ``sign covers''
\(\{R\to R\langle x\rangle,\ R\to R\langle -x\rangle\}\), where
\(R\langle \pm x\rangle\) is the universal monoid under \(R\) in which
\(\pm x\in R^{\gp}\) becomes nonnegative; a compactness argument shows that
every birhizomic cover is a cover for it, and the left Kan extension of
\(F\) turns out to be a sheaf for it with split exact Mayer--Vietoris
sequences, hence acyclic by an induction on the degree.  Acyclicity then
descends to the coarser birhizomic topology.
\end{abstract}

\section{Introduction}

\subsection{Conventions and the objects in the question}
\label{sec:conventions}

Throughout, a \emph{monoid} means a commutative unital monoid \(M\), written
additively, that is
\emph{integral} (cancellative, so that \(M\to M^{\gp}\) is injective, where
\(M^{\gp}\) is the group completion),
\emph{saturated} (if \(g\in M^{\gp}\) and \(ng\in M\) for some integer
\(n\ge1\), then \(g\in M\)), and
\emph{sharp} (the only unit is \(0\)).
All monoids, homomorphisms, pushouts and colimits below are taken in this
category, which we denote by \(\mathbf{Mon}\).  It is a reflective
subcategory of commutative monoids (the reflection integralizes, saturates
and sharpens), so it has all colimits; the pushouts we actually need to
compute are given explicitly in
\eqref{eq:intersection}--\eqref{eq:basechange}.  Basic facts about
integral and saturated monoids, localization and sharpening are collected in
Ogus's book \cite[Ch.~I, \S\S1.3--1.4]{ogus}; the short verifications we need
are included.

Two operations on submonoids of an abelian group \(G\) recur constantly.
For a submonoid \(C\subset G\), the \emph{saturation} of \(C\) in \(G\) is
\(C^{\sat}=\{g\in G:ng\in C\text{ for some }n\ge1\}\),
and the \emph{sharpening} of \(C\) is \(C^\sharp=C/C^\times\), the quotient by
the group of units \(C^\times=C\cap(-C)\).  If \(C\) is saturated in \(G\)
then \(C^\sharp\subset G/C^\times\) is integral, saturated and sharp
(if \(n(g+C^\times)\in C/C^\times\), then \(ng+h\in C\) for some
\(h\in C^\times\), hence \(ng\in C\) and \(g\in C\)).

A monoid \(M\) orders its group completion: for \(x,y\in M^{\gp}\) put
\(x\le y\) if \(y-x\in M\).  This is a translation-invariant partial order,
and \(M\) is called \emph{valuative} if the order is total, i.e.\
\(M\cup(-M)=M^{\gp}\).  We note once and for all that \(M^{\gp}\) is
torsion-free for every monoid \(M\): if \(ng=0\) with \(n\ge1\), saturation
puts \(g\) and \(-g\) in \(M\), and sharpness forces \(g=0\).

Fix a monoid \(P\); it is arbitrary (in particular, no finite generation
is assumed), and it stays fixed throughout.  Let \(A_P\) be the category of
monoids under \(P\),
i.e.\ of homomorphisms \(P\to Q\); we often abbreviate an object to \(Q\),
and we say that \(a\in P\) \emph{dies} in \(Q\), or is \emph{killed} by
\(Q\), if its image in \(Q\) is \(0\).
For \(x\in P\), the \emph{localization} \(P[-x]=P+\N(-x)\subset P^{\gp}\) is
saturated in \(P^{\gp}\) (if \(mg=p-nx\) with \(p\in P\), then
\(m(g+nx)=p+(m-1)nx\in P\), so \(g+nx\in P\) by saturation of \(P\) and
\(g=(g+nx)-nx\in P[-x]\)), and
\[
   L_x:=P[-x]^\sharp
\]
is the initial monoid under \(P\) in which \(x\) dies: a homomorphism
\(P\to T\) sending \(x\) to \(0\) extends uniquely to \(P[-x]\) and then
kills the units of \(P[-x]\) because \(T\) is sharp.  Let \(B_P\subset A_P\)
be the full subcategory with objects \(P\to L_x\), \(x\in P\).  Note
\(L_0=P\).

A \emph{presheaf} on the category \(A_P^{\mathrm{op}}\) is a functor
\((A_P^{\mathrm{op}})^{\mathrm{op}}=A_P\to\mathbf{Set}\) (or \(\to\Ab\)),
hence a \emph{covariant} functor on monoids under \(P\); one should think of
\(A_P^{\mathrm{op}}\) as a category of affine ``monoid schemes''
\(\operatorname{Spec}Q\).  Accordingly, a family of homomorphisms
\(\{R\to Q_i\}\) in \(A_P\) is the same as a family of maps
\(\{Q_i\to R\}\) in \(A_P^{\mathrm{op}}\), and we refer to it as a family
\emph{over} \(R\); fibre products in \(A_P^{\mathrm{op}}\) over \(R\) are
pushouts in \(A_P\) under \(R\).  As usual for such large sites we work in
sufficiently large fixed universes, so that sheafification and injective
resolutions are available.

\subsection{The birhizomic topology}

We recall the definitions from the question.  For a monoid \(M\),
\(x\in M^{\gp}\) and \(y\in M\), write \(x\prec y\) if there are integers
\(m,n\) with \(my\le x\le ny\).  (This boundedness notation also appears in
Molcho--Wise \cite[Def.~2.1.3.1]{MW}.)  A homomorphism \(f:M\to N\) is
\emph{infinitesimal} if, whenever \(x\in M^{\gp}\) and \(y\in M\) satisfy
\(f(x)\in N\), \(x\prec y\) and \(f(y)\prec f(x)\), one has \(x\in M\).
If \(R\xrightarrow{f}Q\xrightarrow{g}N\) are homomorphisms, \(f\) is
\emph{birhizomic at \(g\)} if \(f\) is of finite presentation in
\(\mathbf{Mon}\) (discussed in \S\ref{sec:adjoin}) and for every factorization \(gf=ph\) with
\(R\xrightarrow{h}M\xrightarrow{p}N\) and \(p\) infinitesimal, there is a
unique \(u:Q\to M\) with \(pu=g\) and \(uf=h\).  A family
\(\{R\to Q_i\}_{i\in I}\) is a \emph{birhizomic cover} if for every valuative
monoid \(V\) and every homomorphism \(v:R\to V\) there is a factorization of
\(v\) through some \(Q_i\to V\) such that \(R\to Q_i\) is birhizomic at
\(Q_i\to V\).  The question stipulates that the birhizomic covers (of objects
\(R\) of \(A_P\), all maps being under \(P\)) form a Grothendieck topology on
\(A_P^{\mathrm{op}}\), the \emph{birhizomic topology} \(\Jb\): a sieve on
\(R\) is \(\Jb\)-covering when it contains a birhizomic cover.  (If one
prefers to read \(\Jb\) as the topology \emph{generated} by the birhizomic
covers, nothing below changes: Corollary~\ref{cor:finer} exhibits a
topology \(\tau\) for which every birhizomic cover is covering, and
\(\tau\) then contains \(\Jb\) under either reading.)  These definitions
are taken from the question as stated; a related but different notion of
``infinitesimal'' homomorphism occurs in \cite[Def.~2.1.2.9]{MW}, and we use
only the definition above.

Only two features of this definition enter the proof:
(B1) a map birhizomic at anything is of finite presentation; and
(B2) in a birhizomic cover of \(R\), every homomorphism from \(R\) to
a valuative monoid factors through a member of the family that is of finite
presentation over \(R\).
The infinitesimal lifting condition is never analysed; see
Remark~\ref{rem:infinitesimal}.

\subsection{The functors \texorpdfstring{\(\pi^{-1}\)}{pi\^{}-1} and
  \texorpdfstring{\(\pi_*\)}{pi\_*}, and the main theorem}
\label{sec:functors}

Give \(B_P^{\mathrm{op}}\) the trivial topology, for which every presheaf is a
sheaf, so \(\Sh(B_P^{\mathrm{op}})=\PSh(B_P^{\mathrm{op}})\).  Let
\(j:B_P\hookrightarrow A_P\) be the inclusion.  As in the question,
\(\pi^{-1}:\Sh(B_P^{\mathrm{op}})\to\Sh(A_P^{\mathrm{op}},\Jb)\) is the
composite of the left Kan extension along \(j\) with \(\Jb\)-sheafification,
and \(\pi_*\) is its right adjoint.  Concretely, for a presheaf \(F\) on
\(B_P^{\mathrm{op}}\), i.e.\ a functor \(F:B_P\to\Ab\), the left Kan extension
is the presheaf
\begin{equation}
  (\Lan_jF)(P\to Q)
   =\varinjlim_{(L_a\to Q)\in B_P/Q}F(L_a),
  \label{eq:kan}
\end{equation}
the colimit over the comma category \(B_P/Q\) whose objects are pairs
\((L_a,\ L_a\to Q\text{ under }P)\) \cite[Tags 00VC, 00VE]{stacks}; and
\(\pi_*\) is restriction along \(j\): for a sheaf \(E\) on
\(A_P^{\mathrm{op}}\),
\(\Hom(\pi^{-1}F,E)=\Hom_{\PSh}(\Lan_jF,E)=\Hom(F,E\circ j)\), so
\(\pi_*E=E\circ j\).  The same formulas define \(\pi^{-1}\) and \(\pi_*\) on
abelian (pre)sheaves \cite[Tag 00VF]{stacks}; we check in \S\ref{sec:kan}
that the colimit
\eqref{eq:kan} is filtered, so that the Kan extension of an abelian presheaf
computed in \(\Ab\) agrees with the one computed in sets.  Since \(\pi_*\) is
left exact on abelian sheaves, it has a right derived functor \(\R\pi_*\),
and the unit of the adjunction gives a natural map
\(F\to\pi_*\pi^{-1}F\to\R\pi_*\pi^{-1}F\).

\begin{theorem}\label{thm:main}
Let \(P\) be an arbitrary monoid (integral, saturated and sharp, as
always).  For every presheaf of
abelian groups \(F\) on \(B_P^{\mathrm{op}}\), the unit
\[
             F\longrightarrow \R\pi_*\pi^{-1}F
\]
is a quasi-isomorphism of complexes of abelian presheaves on
\(B_P^{\mathrm{op}}\), i.e.\ an isomorphism in the derived category of
the category \(\Ab\PSh(B_P^{\mathrm{op}})\) of abelian presheaves.  More precisely, the left Kan extension
\(\G=\Lan_jF\) of \eqref{eq:kan} is already a \(\Jb\)-sheaf, so
\(\pi^{-1}F=\G\); the unit \(F\to\pi_*\G\) is an isomorphism; and
\(H^n_{\Jb}(Q,\G)=0\) for every \(n>0\) and every object \(Q\) of \(A_P\),
where \(H^n_{\Jb}(Q,-)\) denotes the \(n\)th right derived functor of
evaluation at \(Q\) on abelian \(\Jb\)-sheaves (\S\ref{sec:acyclic}).
\end{theorem}

The title refers to the following reformulation: since \(\pi^{-1}\) is
exact (the colimit \eqref{eq:kan} is filtered and sheafification is exact),
it is left adjoint to \(\R\pi_*\) between the bounded below derived
categories of \(\Ab\PSh(B_P^{\mathrm{op}})\) and of
\(\Ab\Sh(A_P^{\mathrm{op}},\Jb)\).  Taking \(E=\pi^{-1}F'[n]\) in
\(\Hom(\pi^{-1}F,E)=\Hom(F,\R\pi_*E)\) and using
\(F'\simeq\R\pi_*\pi^{-1}F'\), the theorem gives that the map
\(\operatorname{Ext}^n(F,F')\to\operatorname{Ext}^n(\pi^{-1}F,\pi^{-1}F')\)
induced by \(\pi^{-1}\) (which, by naturality of the unit, is the
composite of these two identifications) is bijective for all \(F,F'\) and
\(n\), the Ext groups being computed in the two categories just named.

\subsection{Strategy}

Cohomology for the birhizomic topology is awkward to compute directly, since
its covers are defined by a lifting property.  We therefore work with an
auxiliary topology \(\tau\) on \(A_P^{\mathrm{op}}\) that has an explicit
basis and is \emph{finer} than \(\Jb\).  Its generating covers are the
\emph{sign covers}
\[
   \{\,R\to R\langle x\rangle,\quad R\to R\langle -x\rangle\,\},
   \qquad x\in R^{\gp},
\]
where \(R\langle x\rangle\) is the universal monoid under \(R\) in which
\(x\) becomes \(\ge0\) (\S\ref{sec:adjoin}).  Geometrically, if \(x=a-b\) with
\(a,b\in R\), these are the two charts of the blow-up of the ideal
\((a,b)\), saturated and sharpened; this chart construction is classical in
toric and logarithmic geometry, see e.g.\ Nizio{\l} \cite[\S4]{niziol}.
After the preliminaries on \(R\langle x\rangle\) in \S\ref{sec:adjoin},
the proof has four steps.

\begin{enumerate}
\item \emph{Comparison of topologies} (\S\ref{sec:tau}).  Every birhizomic
cover is a \(\tau\)-cover (Corollary~\ref{cor:finer}).  By (B2) this reduces
to Proposition~\ref{prop:valuative}: a family of finitely presented maps
\(R\to Q_i\) through which every valuative map \(R\to V\) factors is refined
by finitely many iterated sign covers.  Since \(R\) may be infinitely
generated, there is no a priori bound on the number of sign choices
needed; we obtain one by a compactness argument on the space of ``total
precones'', the saturated submonoids \(D\subset R^{\gp}\) with
\(R\subset D\) and \(D\cup(-D)=R^{\gp}\), a closed subset of
\(\{0,1\}^{R^{\gp}}\).
\item \emph{The Kan extension is a \(\tau\)-sheaf with split Mayer--Vietoris
sequences} (\S\ref{sec:kan}).  The value \(\G(Q)\) depends only on which
elements of \(P\) die in \(Q\), so \(\G\) is insensitive to
\emph{zero-reflecting} maps, i.e.\ homomorphisms \(f\) with
\(f^{-1}(0)=\{0\}\).  Write \(R_\pm=R\langle\pm x\rangle\) for the two
members of a sign cover and \(R_0=R\langle x,-x\rangle\) for their
``intersection'' \(R_+\amalg_RR_-\).  The maps from \(R\) to
\(R_+\), \(R_-\) and \(R_0\) are zero-reflecting unless \(x\) or \(-x\)
already lies in \(R\), in which case the cover contains the identity.  In
both cases the Mayer--Vietoris sequence
\(0\to\G(R)\to\G(R_+)\oplus\G(R_-)\to\G(R_0)\to0\) is split exact
(Lemma~\ref{lem:mvsections}).
\item \emph{Acyclicity for \(\tau\)} (\S\ref{sec:acyclic}).  Let
\(H^n_\tau(R,-)\) denote the \(n\)th right derived functor of evaluation at
\(R\) on abelian \(\tau\)-sheaves.  For a \(\tau\)-sheaf \(\E\) whose
sign-cover Mayer--Vietoris sequences are exact on the right, we show
\(H^n_\tau(R,\E)=0\) for all \(n>0\) and all \(R\)
(Lemma~\ref{lem:acyclic}) by induction on \(n\): a class vanishing on both
members of a sign cover vanishes (Mayer--Vietoris); every class vanishes on
some \(\tau\)-cover (locality of cohomology); and \(\tau\)-covers are
refined by trees of sign covers, along which the vanishing propagates from
the leaves to the root.
\item \emph{Descent to the coarser topology and conclusion}
(\S\ref{sec:acyclic}--\ref{sec:unit}).  Since \(\tau\)-sheafification is
exact, the inclusion of \(\tau\)-sheaves into \(\Jb\)-sheaves preserves
injectives.  A \(\tau\)-injective resolution of
\(\G\) is exact on every object by step~3, hence is also a
\(\Jb\)-injective resolution, so \(H^n_{\Jb}(Q,\G)=0\) for \(n>0\).  Since
\(\pi_*\) is restriction and \(\G\) restricts to \(F\), this is
Theorem~\ref{thm:main}.
\end{enumerate}

\section{Adjoining inequalities}
\label{sec:adjoin}

The basic construction is the universal way of forcing finitely many
elements of \(R^{\gp}\) to become nonnegative.

\begin{definition}\label{def:adjoin}
Let \(R\) be a monoid, \(G=R^{\gp}\), and \(S\subset G\) a subset.  Put
\begin{equation}
 R\langle S\rangle
  =\Bigl(\bigl(R+\textstyle\sum_{s\in S}\N s\bigr)^{\sat}\Bigr)^\sharp ,
 \label{eq:adjoin}
\end{equation}
where the saturation is taken inside \(G\) and the sharpening is as in
\S\ref{sec:conventions}; thus if \(C=(R+\sum_{s\in S}\N s)^{\sat}\), then
\(R\langle S\rangle=C/(C\cap(-C))\), a monoid with group completion
\(G/(C\cap(-C))\).  It comes with the homomorphism
\(R\to R\langle S\rangle\) induced by \(R\subset C\).  For finite sets we
write \(R\langle s_1,\dots,s_k\rangle\).
\end{definition}

\begin{lemma}[Universal property]\label{lem:universal}
For every monoid \(T\), composition with \(R\to R\langle S\rangle\)
identifies
\begin{equation}
 \Hom(R\langle S\rangle,T)
 =
 \{\,a:R\longrightarrow T:
                    a^{\gp}(s)\in T\text{ for every }s\in S\,\},
 \label{eq:universal}
\end{equation}
where \(a^{\gp}:G\to T^{\gp}\) is the induced homomorphism of groups.  In
particular \(R\to R\langle S\rangle\) is an epimorphism in \(\mathbf{Mon}\),
and if \(S\) is finite it is of finite presentation.
\end{lemma}

\begin{proof}
Let \(a:R\to T\) with \(a^{\gp}(S)\subset T\).  Then \(a^{\gp}\) maps
\(R+\sum_{s\in S}\N s\) into \(T\); it maps \(C\) into \(T\) because \(T\) is
saturated (if \(ng\in R+\sum\N s\) then \(na^{\gp}(g)\in T\), so
\(a^{\gp}(g)\in T\)); and it kills \(C\cap(-C)\) because \(T\) is sharp.
Hence \(a^{\gp}\) induces \(C/(C\cap(-C))\to T\) extending \(a\).  The
extension is unique, because a homomorphism out of the integral monoid
\(R\langle S\rangle\) is determined by the induced map on group completions,
and \(G\to G/(C\cap(-C))\) is surjective.  Conversely, for any
\(R\langle S\rangle\to T\) the composite \(a\) with \(R\to R\langle S\rangle\)
has \(a^{\gp}(s)\in T\), since \(s\) maps into \(R\langle S\rangle\).
Uniqueness of the extension says exactly that \(R\to R\langle S\rangle\) is an
epimorphism.

For finite presentation, write each \(s\in S\) as \(s=p_s-q_s\) with
\(p_s,q_s\in R\).  Then \eqref{eq:universal} says that \(R\langle S\rangle\)
is obtained from \(R\) by adjoining one generator \(t_s\) and imposing the
relation \(q_s+t_s=p_s\), for each \(s\in S\), in \(\mathbf{Mon}\): a
homomorphism \(R\langle S\rangle\to T\) is a homomorphism \(a:R\to T\)
together with elements \(t_s\in T\) satisfying \(a(q_s)+t_s=a(p_s)\), the
\(t_s=a^{\gp}(s)\) being unique by integrality.  In other words,
\(R\langle S\rangle\) is the quotient, in \(\mathbf{Mon}\), of
\(R\oplus\N^S\) by the finitely many relations \(q_s+t_s=p_s\).
\end{proof}

The question speaks of finite presentation ``in the category of (integral,
saturated, sharp) monoids''.  One may read this as a presentation by
finitely many generators and relations in \(\mathbf{Mon}\) (\(Q\) is the
quotient in \(\mathbf{Mon}\) of \(R\oplus\N^n\) by \(m\) relations
\(u_i=v_i\) with \(u_i,v_i\in R\oplus\N^n\), so that \(\Hom_R(Q,T)\), the
set of homomorphisms under \(R\), is the set of \(t\in T^n\) satisfying the \(m\) equations obtained by
evaluating \(u_i,v_i\) at \(a:R\to T\) and \(t\)), or categorically.  Under
either reading, the property we shall use is the usual one: \emph{for
\(Q\) of finite presentation over \(R\), the functor \(\Hom_R(Q,-)\) on
monoids under \(R\) commutes with filtered colimits}.  The categorical
reading is this property by definition; for the presentation reading it
follows from the next lemma, because \(\Hom_R(Q,T)\) is then the equalizer
of two maps \(T^n\rightrightarrows T^m\) natural in \(T\), a finite limit
of functors \(T\mapsto T^k\), which commute with filtered colimits computed
on underlying sets.

\begin{lemma}\label{lem:filtered}
Filtered colimits exist in \(\mathbf{Mon}\) and are computed as filtered
colimits of ordinary commutative monoids, i.e.\ on underlying sets.
\end{lemma}

\begin{proof}
Let \(T=\varinjlim T_\alpha\) be a filtered colimit of ordinary commutative
monoids, with \(T_\alpha\) integral, saturated and sharp.  Every element of
\(T\) comes from some stage and every equality in \(T\) holds at some later
stage; so cancellation (\(u+v=u+w\Rightarrow v=w\)) and sharpness
(\(u+v=0\Rightarrow u=v=0\)) pass from the \(T_\alpha\) to \(T\).
Saturation: if \(g=u-v\in T^{\gp}\) with \(u,v\in T\) and \(ng\in T\), say
\(nu=nv+w\) with \(w\in T\), lift \(u,v,w\) and this equality to a stage
\(T_\alpha\); there \(n(u-v)=w\in T_\alpha\), so \(u-v\in T_\alpha\) by
saturation, and its image in \(T\) is \(g\).  Thus \(T\) lies in
\(\mathbf{Mon}\), and it is then the colimit there.  (Integrality and
saturation are also covered by \cite[Prop.~I.1.3.6]{ogus}.)
\end{proof}

Two formal consequences of \eqref{eq:universal}, both proved by comparing
the functors \(\Hom(-,T)\), are used repeatedly: for subsets
\(S,S'\subset G\) and a homomorphism \(u:R\to R'\),
\begin{align}
 R\langle S\rangle\amalg_R R\langle S'\rangle
       &\simeq R\langle S\cup S'\rangle, \label{eq:intersection}\\
 R'\amalg_R R\langle S\rangle
       &\simeq R'\langle u^{\gp}(S)\rangle .\label{eq:basechange}
\end{align}
Indeed, a map out of the left side of \eqref{eq:intersection} to \(T\) is a
map \(a:R\to T\) with \(a^{\gp}(S)\subset T\) and \(a^{\gp}(S')\subset T\),
and a map out of the left side of \eqref{eq:basechange} is a map
\(b:R'\to T\) with \((bu)^{\gp}(S)=b^{\gp}(u^{\gp}(S))\subset T\).
In particular \(R\langle S\rangle\langle u^{\gp}(S')\rangle
\simeq R\langle S\cup S'\rangle\) for \(u:R\to R\langle S\rangle\).

The remaining two lemmas of this section are the algebraic heart of the
argument: the first produces valuative monoids, the second says that
adjoining a sign is harmless from the point of view of which elements become
zero.

\begin{lemma}[Extension of an order]\label{lem:order}
Every monoid \(R\) embeds into a valuative monoid \(W\) with
\(W^{\gp}=R^{\gp}\).
\end{lemma}

This is \cite[Lemma~2.1.2.11]{MW}; we include the short proof.

\begin{proof}
Let \(G=R^{\gp}\).  Saturated sharp submonoids of \(G\) containing \(R\) form
a poset under inclusion in which every chain has an upper bound, its union
(saturation and sharpness are conditions on finitely many elements at a
time).  By Zorn's lemma choose a maximal one, \(C\).  We claim
\(C\cup(-C)=G\), so that \(W=C\) is valuative with \(W^{\gp}=G\).

Suppose that neither \(x\) nor \(-x\) belongs to \(C\).  We show that
\(C'=(C+\N x)^{\sat}\) is still sharp; since \(C'\) is saturated, contains
\(R\), and contains \(x\notin C\), this contradicts the maximality of \(C\).  Let \(u,-u\in C'\).  Then
there are \(m,n>0\), \(a,b\in C\) and \(i,j\geq0\) such that
\[
       mu=a+ix,\qquad -nu=b+jx .
\]
Multiplying the first equation by \(n\), the second by \(m\), and adding,
\[
       0=na+mb+(ni+mj)x.
\]
If \(ni+mj>0\), then \(-(ni+mj)x=na+mb\in C\), and saturation of \(C\) gives
\(-x\in C\), a contradiction.  Otherwise \(i=j=0\), so \(mu\) and \(-nu\) lie
in \(C\); then \(mnu\) and \(-mnu\) lie in \(C\), sharpness of \(C\) gives
\(mnu=0\), and torsion-freeness of \(G\) gives \(u=0\).
\end{proof}

Call a homomorphism of monoids \(f:Q\to Q'\) \emph{zero-reflecting} if
\(f^{-1}(0)=\{0\}\).  (Between sharp monoids this is the condition that
\(f\) reflects units, called \emph{local} in \cite[p.~15]{ogus}.)

\begin{lemma}\label{lem:zero}
Let \(x\in R^{\gp}\), and set
\[
 R_+=R\langle x\rangle,\qquad
 R_-=R\langle -x\rangle,\qquad
 R_0=R\langle x,-x\rangle .
\]
\begin{enumerate}
\item If neither \(x\) nor \(-x\) lies in \(R\), the three maps
\(R\to R_+\), \(R\to R_-\) and \(R\to R_0\) are zero-reflecting.
\item If \(x\in R\), then \(R_+=R\), and \(R_0=R_-\) with \(R_-\to R_0\) the
identity and \(R_+\to R_0\) the map \(R\to R_-\).  Symmetrically, if
\(-x\in R\), then \(R_-=R\) and \(R_0=R_+\).
\end{enumerate}
\end{lemma}

\begin{proof}
(2) If \(x\in R\), then \(R+\N x=R\), which is saturated and sharp, so
\(R_+=R\); and \(R\langle x,-x\rangle=R\langle -x\rangle\) by
\eqref{eq:adjoin}, since \(\N x\subset R\).  The description of the maps
follows from the uniqueness in Lemma~\ref{lem:universal}.  The case
\(-x\in R\) is obtained by replacing \(x\) with \(-x\).

(1) Since \(R\to R_0\) factors through \(R\to R_+\) and through \(R\to R_-\)
(Lemma~\ref{lem:universal}), it suffices to treat \(R_0\).  Before
sharpening, \(R_0\) is \(C=(R+\Z x)^{\sat}\).  If \(r\in R\) maps to zero in
\(C^\sharp\), then \(-r\in C\), so
\[
                   -nr=a+kx
       \qquad(n>0,\ a\in R,\ k\in\Z).
\]
If \(k>0\), then \(-kx=a+nr\in R\), and saturation gives \(-x\in R\);
if \(k<0\), then \((-k)x=a+nr\in R\) gives \(x\in R\).  Both are excluded.
If \(k=0\), then \(nr\) and \(-nr\) lie in \(R\), so \(nr=0\) by sharpness
and \(r=0\) by torsion-freeness.
\end{proof}

\section{A topology finer than the birhizomic topology}
\label{sec:tau}

\begin{definition}\label{def:tau}
For an object \(R\) of \(A_P\) and \(x\in R^{\gp}\), the \emph{sign cover}
of \(R\) at \(x\) is the family
\begin{equation}
             \{\,R\longrightarrow R\langle x\rangle,\quad
                    R\longrightarrow R\langle -x\rangle\,\}
             \qquad(\text{written in the monoid direction}).
                                                         \label{eq:binary}
\end{equation}
Let \(\tau\) be the Grothendieck topology on \(A_P^{\mathrm{op}}\) generated
by the sign covers, i.e.\ the smallest topology for which the sieve generated
by every sign cover is covering.
\end{definition}

By \eqref{eq:basechange}, the base change of the sign cover at \(x\) along
\(u:R\to R'\) (a fibre product in \(A_P^{\mathrm{op}}\)) is the sign cover of
\(R'\) at \(u^{\gp}(x)\): \emph{sign covers are stable under base change}.
This gives \(\tau\) an explicit description.  A \emph{tree family} over
\(R\) is the family of leaves \(\{R\to T_\ell\}_{\ell\in\Lambda}\) of a finite
binary tree with root \(R\) in which the two children of each internal
vertex \(M\) form a sign cover of \(M\) (\(R\to T_\ell\) being the
composite along the path from the root); its \emph{depth} is the depth of
the tree, the trivial tree (depth \(0\)) giving \(\{\mathrm{id}_R\}\).
Base change along \(u:R\to R'\) takes a tree family over \(R\) to a tree
family over \(R'\) of the same depth (apply \eqref{eq:basechange} vertex
by vertex), and replacing each leaf of a tree family by a tree family over
that leaf (grafting) gives again a tree family.  Consequently the sieves
containing a tree family form a Grothendieck topology (the three axioms---maximal sieve, stability under base change, local
character, i.e.\ transitivity---follow from the trivial tree, base change
and grafting respectively), in
which sign covers are covering; conversely, in any topology in which sign
covers are covering, every tree family generates a covering sieve, by
induction on the depth using the local character axiom.  Hence
\begin{equation}
  \text{a sieve on \(R\) is \(\tau\)-covering if and only if it contains
  a tree family over \(R\).}\label{eq:treesieve}
\end{equation}
For \(\tau\), or any other topology \(J\) on \(A_P^{\mathrm{op}}\), we say
that a family over \(R\) is a \emph{\(J\)-cover} if the sieve it generates
is \(J\)-covering.  Thus a family over \(R\) is a \(\tau\)-cover as soon as
it is \emph{refined} by a tree family, meaning that every leaf \(T_\ell\)
receives a map under \(R\) from some member of the family.

Members of tree families are epimorphisms, being composites of the
epimorphisms \(M\to M\langle\pm x\rangle\) (Lemma~\ref{lem:universal}).
Recall that a presheaf \(\mathcal F\) on \(A_P^{\mathrm{op}}\) satisfies the
\emph{sheaf condition} for a family \(\{R\to T_\ell\}_{\ell\in\Lambda}\) if
\(\mathcal F(R)\to\prod_\ell\mathcal F(T_\ell)\) is injective with image
the \emph{compatible} families \((s_\ell)\), those for which \(s_\ell\) and
\(s_m\) have the same image in \(\mathcal F(T_\ell\amalg_RT_m)\) for all
\(\ell,m\) (fibre products in \(A_P^{\mathrm{op}}\) being pushouts in
\(A_P\)); when \(R\to T_\ell\) is an epimorphism,
\(T_\ell\amalg_RT_\ell=T_\ell\) and the condition for \(\ell=m\) is vacuous.
For a sign cover, compatibility of \((s_+,s_-)\) means that \(s_+\) and
\(s_-\) have the same image in
\(\mathcal F(R\langle x\rangle\amalg_RR\langle -x\rangle)
=\mathcal F(R\langle x,-x\rangle)\) \eqref{eq:intersection}.

\begin{lemma}[Sheaf criterion]\label{lem:criterion}
A presheaf \(\mathcal F\) of sets or of abelian groups on
\(A_P^{\mathrm{op}}\) is a \(\tau\)-sheaf as soon as it satisfies the sheaf
condition for every sign cover \eqref{eq:binary}.
\end{lemma}

\begin{proof}
\emph{Step A: \(\mathcal F\) satisfies the sheaf condition for every tree
family, over every object.}  We induct on the depth \(d\); depth \(0\) is
trivial and depth \(1\) is the hypothesis.  Let \(\{R\to T_\ell\}\) have
depth \(d+1\ge2\), let \(R_\pm\) be the children of the root and
\(\Lambda=\Lambda_+\sqcup\Lambda_-\) the leaves of the two subtrees, which
are tree families over \(R_\pm\) of depth \(\le d\).  Let \((s_\ell)\) be
compatible.  For \(\ell,m\in\Lambda_+\), restricting along the canonical map
\(T_\ell\amalg_RT_m\to T_\ell\amalg_{R_+}T_m\) shows that
\((s_\ell)_{\ell\in\Lambda_+}\) is compatible over \(R_+\), so by induction
there is a unique \(s_+\in\mathcal F(R_+)\) restricting to \(s_\ell\) for
all \(\ell\in\Lambda_+\); similarly \(s_-\).  We claim that \(s_+\) and
\(s_-\) have the same image in \(\mathcal F(R_0)\), \(R_0=R_+\amalg_RR_-\).
The base change of the \(R_+\)-subtree along \(R_+\to R_0\) is a tree
family over \(R_0\) of depth \(\le d\) with leaves
\(T_\ell\amalg_{R_+}R_0=T_\ell\amalg_RR_-\) (\(\ell\in\Lambda_+\)), so by
the injectivity part of the induction hypothesis it suffices to compare the
restrictions of \(s_+\) and \(s_-\) to each \(T_\ell\amalg_RR_-\); these are
the restrictions of \(s_\ell\) and of \(s_-\).  In turn, the base change of
the \(R_-\)-subtree along \(R_-\to T_\ell\amalg_RR_-\) is a tree family of
depth \(\le d\) with leaves \(T_\ell\amalg_RT_m\) (\(m\in\Lambda_-\)), on
which \(s_\ell\) and \(s_-\) restrict to the images of \(s_\ell\) and
\(s_m\), which agree by compatibility; injectivity once more proves the
claim.  By the sheaf condition for the sign cover \(\{R\to R_\pm\}\) there
is a unique \(s\in\mathcal F(R)\) restricting to \(s_\pm\), hence to all
\(s_\ell\).  If \(s'\in\mathcal F(R)\) also restricts to all \(s_\ell\), then
\(s'|_{R_\pm}=s_\pm\) by the uniqueness in the induction hypothesis, so
\(s'=s\).

\emph{Step B: \(\mathcal F\) is a \(\tau\)-sheaf} (this is the usual
comparison between sheaves for a site and for its associated topology,
cf.\ \cite[Tag 00ZC]{stacks}; we include the argument since tree families,
not being closed under isomorphism, do not form a site in the sense of that
tag).  Let \(\Sigma\) be a
\(\tau\)-covering sieve on \(R\); by \eqref{eq:treesieve} it contains a
tree family \(\{R\to T_\ell\}\).  A compatible family on \(\Sigma\) is a
collection of sections \(s_V\in\mathcal F(V)\), one for each arrow
\((R\to V)\in\Sigma\) (sections are indexed by arrows; we write \(s_V\)
when the arrow is clear), with \(s_{V'}=s_V|_{V'}\) for every \(V\to V'\)
under \(R\), and we must show that \(s\mapsto(s|_V)_V\) is a bijection from
\(\mathcal F(R)\) onto such collections.  Injectivity follows from Step~A
for \(\{R\to T_\ell\}\).  Given \((s_V)\), the sections \(s_{T_\ell}\) are
compatible (both \(s_{T_\ell}\) and \(s_{T_m}\) restrict to
\(s_{T_\ell\amalg_RT_m}\)), so Step~A gives \(s\in\mathcal F(R)\) with
\(s|_{T_\ell}=s_{T_\ell}\) for all \(\ell\).  For \((R\to V)\in\Sigma\), the
sections \(s|_V\) and \(s_V\) both restrict to \(s_{V\amalg_RT_\ell}\) on
each leaf of the base-changed tree family \(\{V\to V\amalg_RT_\ell\}\),
hence are equal by Step~A.
\end{proof}

Our goal in this section is to show that \(\tau\) is finer than \(\Jb\).
The next proposition concerns finitely presented families that ``see'' all
valuative monoids; by (B2), it applies to birhizomic covers.  It contains
the compactness argument needed for possibly infinitely generated \(R\).

\begin{proposition}\label{prop:valuative}
Let \(R\) be an object of \(A_P\) and let \(\{R\to Q_i\}_{i\in I}\) be a
family of maps of finite presentation such that every homomorphism
\(R\to V\) to a valuative monoid \(V\) factors through some \(R\to Q_i\).
Then the family is refined by a tree family; in particular it is a
\(\tau\)-cover.
\end{proposition}

\begin{proof}
\emph{Step 1: a compact space of orders.}
Let \(G=R^{\gp}\), and let \(\V(R)\) be the set of saturated submonoids
\(D\subset G\) with
\[
                         R\subset D,\qquad D\cup(-D)=G .
\]
Such a \(D\) is a ``total precone'': it induces a total preorder on \(G\),
but it need not be sharp, and it is essential to allow this: the precones
\(D_\epsilon\) produced in Step~5 below contain the kernel of a surjection
\(G\to T_\epsilon^{\gp}\) and are not sharp in general.  Identify \(D\) with its characteristic
function and regard \(\V(R)\) as a subset of the product \(\{0,1\}^{G}\),
compact by Tychonoff's theorem.  It is closed: each of the conditions
\(r\in D\) (for \(r\in R\)), \(0\in D\), ``\(g,h\in D\Rightarrow g+h\in D\)'',
``\(ng\in D\Rightarrow g\in D\)'' and ``\(g\in D\) or \(-g\in D\)'' involves
finitely many coordinates, and \(\V(R)\) is the intersection of these closed
conditions.  Hence \(\V(R)\) is compact.  (The same device, embedding a set of
valuation-like objects into a Boolean product, appears in
\cite[Lemma~5.3.6]{CT}.)  For finite \(S\subset G\), the subset
\[
                       U_S=\{D\in\V(R):S\subset D\}
\]
is open and closed, being defined by finitely many coordinates.

\emph{Step 2: the valuative monoid of a precone as a filtered colimit.}
For \(D\in\V(R)\), write \(H_D=D\cap(-D)\) and \(V_D=D/H_D=D^\sharp\).
Then \(V_D\) is a monoid (integral, saturated, sharp, cf.\
\S\ref{sec:conventions}), valuative because \(D\cup(-D)=G\), and a monoid
under \(R\) via \(R\subset D\to V_D\).  We claim that, in the category of
monoids under \(R\),
\begin{equation}
       V_D=\varinjlim_{\substack{S\subset D\\ S\ {\rm finite}}}
                                         R\langle S\rangle,
                         \label{eq:valcolim}
\end{equation}
the colimit being over the directed set of finite subsets of \(D\), with
transition maps \(R\langle S\rangle\to R\langle S'\rangle\) for
\(S\subset S'\) given by Lemma~\ref{lem:universal}.  Indeed, by
\eqref{eq:universal} a compatible family of maps \(R\langle S\rangle\to T\)
under \(a:R\to T\) is the same as the condition \(a^{\gp}(D)\subset T\).
Given this condition, \(a^{\gp}|_D:D\to T\) kills \(H_D\) (as \(T\) is sharp)
and so factors uniquely through \(V_D\); conversely a map \(V_D\to T\) under
\(R\) composed with \(D\to V_D\) shows \(a^{\gp}(D)\subset T\).  So \(V_D\)
represents the functor \(\varprojlim_S\Hom_R(R\langle S\rangle,-)\), which is
the claim.

\emph{Step 3: finite presentation yields finite stages.}
Fix \(D\in\V(R)\).  By hypothesis the valuative map \(R\to V_D\) factors as
\(R\to Q_i\to V_D\) for some \(i\).  Since \eqref{eq:valcolim} is a
filtered colimit of monoids under \(R\) and \(\Hom_R(Q_i,-)\) commutes with
such colimits (finite presentation; see the discussion following
Lemma~\ref{lem:universal}), the map \(Q_i\to V_D\) factors as
\[
                       Q_i\longrightarrow R\langle S_D\rangle
                                      \longrightarrow V_D
\]
under \(R\), for some finite \(S_D\subset D\).  The clopen sets \(U_{S_D}\),
\(D\in\V(R)\), cover \(\V(R)\), since \(D\in U_{S_D}\).  By compactness choose
\(D_1,\dots,D_m\) with \(\V(R)=U_{S_1}\cup\cdots\cup U_{S_m}\), where
\(S_j=S_{D_j}\), and retain the maps
\begin{equation}
                         Q_{i_j}\longrightarrow R\langle S_j\rangle
                         \qquad(j=1,\dots,m)
                         \label{eq:retained}
\end{equation}
under \(R\).  Put \(S=S_1\cup\cdots\cup S_m\), a finite subset of \(G\).

\emph{Step 4: the sign tree.}
Enumerate \(S=\{s_1,\dots,s_k\}\) and form the binary tree of depth \(k\)
whose vertex at level \(\ell\) with sign history
\(\epsilon(s_1),\dots,\epsilon(s_\ell)\) is
\(R\langle\epsilon(s_1)s_1,\dots,\epsilon(s_\ell)s_\ell\rangle\); by
\eqref{eq:basechange} and \eqref{eq:intersection}, the sign cover of this
vertex at the image of \(s_{\ell+1}\) has as members the two vertices at the
next level.  The leaves are
\begin{equation}
      T_\epsilon
          =R\langle\,\epsilon(s)s:s\in S\,\rangle,
             \qquad \epsilon:S\longrightarrow\{1,-1\},
             \label{eq:signcover}
\end{equation}
so \(\{R\to T_\epsilon\}_\epsilon\) is a tree family.  It remains to show
that each \(R\to T_\epsilon\) factors through some
\(R\langle S_j\rangle\); composing with \eqref{eq:retained} then shows that
the tree family refines \(\{R\to Q_i\}\).

\emph{Step 5: each leaf factors through some \(R\langle S_j\rangle\).}
The idea is to attach to each leaf \(T_\epsilon\) a precone
\(D_\epsilon\in\V(R)\) extending the inequalities \(\epsilon(s)s\ge0\)
that define \(T_\epsilon\); the finite subcover gives
\(S_j\subset D_\epsilon\) for some \(j\), and we then check that the
elements of \(S_j\) are nonnegative already in \(T_\epsilon\).
Fix \(\epsilon\).  By Lemma~\ref{lem:order}, embed \(T_\epsilon\) into a
valuative monoid \(W_\epsilon\) with
\(W_\epsilon^{\gp}=T_\epsilon^{\gp}\).  Let \(D_\epsilon\subset G\) be the
inverse image of \(W_\epsilon\) under the group homomorphism
\(G\to T_\epsilon^{\gp}=W_\epsilon^{\gp}\).  Then \(D_\epsilon\in\V(R)\):
it contains \(R\) (which maps into \(T_\epsilon\subset W_\epsilon\)), it is
saturated (as the inverse image of a saturated monoid under a group
homomorphism), and \(D_\epsilon\cup(-D_\epsilon)=G\) because
\(W_\epsilon\cup(-W_\epsilon)=W_\epsilon^{\gp}\).  Choose \(j\) with
\(D_\epsilon\in U_{S_j}\), i.e.\ \(S_j\subset D_\epsilon\).

We check that every \(s\in S_j\) becomes nonnegative in \(T_\epsilon\).
Let \(\bar s\in T_\epsilon^{\gp}\) be the image of \(s\).  If
\(\epsilon(s)=1\), then \(\bar s\in T_\epsilon\) by construction of
\(T_\epsilon\).  If \(\epsilon(s)=-1\), then \(-\bar s\in T_\epsilon\subset
W_\epsilon\), while \(s\in D_\epsilon\) says \(\bar s\in W_\epsilon\);
sharpness of \(W_\epsilon\) gives \(\bar s=0\) in
\(W_\epsilon^{\gp}=T_\epsilon^{\gp}\), so again \(\bar s\in T_\epsilon\).
(Here it matters that \(W_\epsilon\) has the \emph{same} group completion as
\(T_\epsilon\): a map to an arbitrary valuative monoid would only give
\(\bar s\mapsto0\) there.)  By Lemma~\ref{lem:universal} there is a map
\(R\langle S_j\rangle\to T_\epsilon\) under \(R\), as required.
\end{proof}

\begin{corollary}\label{cor:finer}
Every birhizomic covering sieve on \(A_P^{\mathrm{op}}\) is \(\tau\)-covering;
that is, \(\Jb\) is coarser than \(\tau\), \(\Jb\subset\tau\).  Consequently
every \(\tau\)-sheaf is a \(\Jb\)-sheaf.
\end{corollary}

\begin{proof}
Let \(\{R\to Q_i\}_{i\in I}\) be a birhizomic cover and let \(I'\subset I\)
be the set of indices for which \(R\to Q_i\) is of finite presentation.
By (B1) and the definition of a birhizomic cover, every homomorphism from
\(R\) to a valuative monoid factors through some \(Q_i\) with \(i\in I'\).
Proposition~\ref{prop:valuative} applied to \(\{R\to Q_i\}_{i\in I'}\) shows
that the sieve generated by this subfamily, hence also the sieve generated
by the whole family, is \(\tau\)-covering.  Every \(\Jb\)-covering sieve
contains a birhizomic cover and is therefore \(\tau\)-covering.  The last
assertion follows, since a \(\tau\)-sheaf satisfies the sheaf condition for
every \(\tau\)-covering sieve, in particular for every \(\Jb\)-covering one.
\end{proof}

\section{The Kan extension and its Mayer--Vietoris sequences}
\label{sec:kan}

Fix a presheaf of abelian groups \(F\) on \(B_P^{\mathrm{op}}\), i.e.\ a
functor \(F:B_P\to\Ab\), and let \(\G=\Lan_jF\) be its left Kan extension
\eqref{eq:kan} to \(A_P\), a presheaf of abelian groups on
\(A_P^{\mathrm{op}}\):
\[
       \G(P\to Q)=
              \varinjlim_{(L_a\to Q)\in B_P/Q} F(L_a).
\]
We first make the indexing category explicit.  By the universal property of
\(L_a=P[-a]^\sharp\), a map \(L_a\to Q\) under \(P\) exists if and only if
the image of \(a\) in \(Q\) is zero, and it is then unique.  In particular
\(B_P\) is a \emph{thin} category (a preorder): there is at most one map
\(L_a\to L_b\), and it exists if and only if \(a\mapsto0\) in \(L_b\).  So
\(B_P/Q\) is the preorder of elements \(a\in P\) that die in \(Q\)
(elements giving the same object of \(B_P\) being identified), with \(a\le b\) when
\(a\) dies in \(L_b\).  It is nonempty, since it contains \(L_0=P\),
and directed: if \(a\) and \(b\) die in \(Q\), so does \(a+b\), and
\(a,b\mapsto0\) in \(L_{a+b}\) because the images of \(a\) and \(b\) in the
sharp monoid \(L_{a+b}\) sum to zero.  Thus \eqref{eq:kan} is a filtered
colimit, and since filtered colimits of abelian groups are computed on
underlying sets, the Kan extension of \(F\) in abelian presheaves is the same
as in presheaves of sets, as announced in the introduction.

Two consequences will be used.

\begin{lemma}\label{lem:kanprops}
\begin{enumerate}
\item For \(Q=L_a\) in \(B_P\), the comma category \(B_P/L_a\) has the
terminal object \((L_a,\mathrm{id})\), so the unit of the Kan extension is an
isomorphism
\begin{equation}
                              F(L_a)\xrightarrow{\ \sim\ }\G(L_a).
                     \label{eq:restrict}
\end{equation}
\item If \(Q\to Q'\) is a zero-reflecting map under \(P\), then an element
\(a\in P\) dies in \(Q\) if and only if it dies in \(Q'\), so the comma
categories \(B_P/Q\) and \(B_P/Q'\) coincide, and the restriction map
\begin{equation}
                              \G(Q)\xrightarrow{\ \sim\ }\G(Q')
                     \label{eq:zeroreflect}
\end{equation}
is an isomorphism (it is the identity of the common colimit).
\end{enumerate}
\end{lemma}

\begin{proof}
(1) Every object \((L_b,L_b\to L_a)\) of \(B_P/L_a\) has the given map as
its unique map to \((L_a,\mathrm{id})\), by thinness; a colimit over a
category with a terminal object is the value at that object, and the unit
is the corresponding structure map.  (2) An element killed in \(Q\) is
killed in \(Q'\); conversely, if \(a\mapsto0\) in \(Q'\), then the image of
\(a\) in \(Q\) is killed by the zero-reflecting map \(Q\to Q'\), so it is
\(0\).  The restriction map \(\G(Q)\to\G(Q')\) is induced by the identity of
the diagrams, hence is the identity.
\end{proof}

In words, \(\G(Q)\) depends only on the set \(\{a\in P:a\text{ dies in }Q\}\)
of elements of \(P\) killed by \(Q\) (a \emph{face} of \(P\): a submonoid
containing \(a\) and \(b\) whenever it contains \(a+b\)), and \(\G\) is
``locally constant'' along zero-reflecting maps.  Combined with Lemma~\ref{lem:zero}, which says that
sign covers consist of zero-reflecting maps except in the degenerate case,
this gives the sheaf property and more.

\begin{lemma}\label{lem:mvsections}
For every sign cover \eqref{eq:binary}, with \(R_\pm\) and \(R_0\) as in
Lemma~\ref{lem:zero}, the sequence
\begin{equation}
 0\longrightarrow\G(R)
   \longrightarrow\G(R_+)\oplus\G(R_-)
   \longrightarrow\G(R_0)\longrightarrow0,
                         \label{eq:mvsections}
\end{equation}
whose first map is the pair of restrictions and whose second map is the
difference of the two restrictions to \(R_0\), is exact, and in fact split.
Moreover \(\G\) is a
\(\tau\)-sheaf, hence (Corollary~\ref{cor:finer}) a \(\Jb\)-sheaf.
\end{lemma}

\begin{proof}
Suppose first that neither \(x\) nor \(-x\) lies in \(R\).  By
Lemma~\ref{lem:zero}(1) the maps \(R\to R_+\), \(R\to R_-\) and
\(R\to R_0\) are zero-reflecting, so by \eqref{eq:zeroreflect} the three
restriction maps out of \(\G(R)\) are isomorphisms; since
\(\G(R)\to\G(R_\pm)\to\G(R_0)\) composes to the isomorphism
\(\G(R)\to\G(R_0)\), the maps \(\G(R_\pm)\to\G(R_0)\) are isomorphisms too.
Identifying all four groups with \(A=\G(R)\) through these isomorphisms,
every restriction map becomes the identity of \(A\) and the sequence
becomes
\[
                   0\longrightarrow A
                    \xrightarrow{a\mapsto(a,a)}A\oplus A
                    \xrightarrow{(a,b)\mapsto a-b}A
                    \longrightarrow0,
\]
which is exact, and split by \(a\mapsto(a,0)\).

If \(x\in R\), then by Lemma~\ref{lem:zero}(2) \(R_+=R\), \(R_0=R_-\), the
map \(R_-\to R_0\) is the identity and \(R_+\to R_0\) is \(R\to R_-\).
Writing \(A=\G(R)\), \(B=\G(R_-)\) and \(\rho:A\to B\) for the restriction,
the sequence reads
\[
                   0\longrightarrow A
                    \xrightarrow{a\mapsto(a,\rho(a))}A\oplus B
                    \xrightarrow{(a,b)\mapsto \rho(a)-b}B
                    \longrightarrow0,
\]
which is exact, split by \((a,b)\mapsto a\), for any homomorphism \(\rho\).
The case \(-x\in R\) is the same with the roles of \(R_+\) and \(R_-\)
exchanged.

For the sheaf property, by Lemma~\ref{lem:criterion} it suffices to check
the sheaf condition for sign covers, and as noted before that lemma this
condition says that \(\G(R)\) is the equalizer of the two restriction maps
\(\G(R_+)\times\G(R_-)\rightrightarrows\G(R_0)\), which is exactness of
\eqref{eq:mvsections} at the first two terms.
\end{proof}

\section{Acyclicity, and passage to the coarser topology}
\label{sec:acyclic}

For a Grothendieck topology \(J\) on \(A_P^{\mathrm{op}}\), an object \(R\)
and an abelian \(J\)-sheaf \(\E\), we write \(H^n_J(R,\E)\) for the \(n\)th
right derived functor of \(\E\mapsto\E(R)\) on the abelian category
\(\Ab\Sh(A_P^{\mathrm{op}},J)\) of abelian \(J\)-sheaves, computed by
injective resolutions.  We shall
use two standard facts.  First, the Mayer--Vietoris sequence
\cite[Tag 0EVY]{stacks} for a cover \(\{R_\pm\to R\}\) in
\(A_P^{\mathrm{op}}\) by two monomorphisms (\(R\to R_\pm\) being
epimorphisms in \(A_P\)) with fibre product \(R_0\); we spell out the
version we need in the proof below.  Second,
\emph{locality of cohomology} \cite[Tag 03P1]{stacks}: for \(n>0\), every
class \(\xi\in H^n_J(R,\E)\) vanishes on the members of some \(J\)-cover
of \(R\) (in the sense of \S\ref{sec:tau}).  (Explicitly: if \(\E\to\I^\bullet\) is an injective
resolution and \(z\in\I^n(R)\) a cocycle representing \(\xi\), the exactness
of \(\I^\bullet\) as a complex of \emph{sheaves} in degree \(n>0\) says that
\(z\) is locally a coboundary.)

\begin{lemma}\label{lem:acyclic}
Let \(\E\) be an abelian \(\tau\)-sheaf on \(A_P^{\mathrm{op}}\) such that for
every sign cover \eqref{eq:binary}, with \(R_\pm\) and \(R_0\) as in
Lemma~\ref{lem:zero}, the difference map
\begin{equation}
                   \E(R_+)\oplus\E(R_-)\longrightarrow\E(R_0)\label{eq:surj}
\end{equation}
is surjective.  Then \(H^n_\tau(R,\E)=0\) for all \(n>0\) and all objects
\(R\) of \(A_P\).
\end{lemma}

\begin{proof}
\emph{Mayer--Vietoris.}  For an object \(Y\) of \(A_P\) let \(h_Y\) be the
presheaf represented by \(Y\) on \(A_P^{\mathrm{op}}\), so
\(h_Y(X)=\Hom_{A_P}(Y,X)\), and let \(\Z_\tau[Y]\) be the \(\tau\)-sheaf
associated with the free abelian presheaf \(\Z[h_Y]\); then
\(\Hom(\Z_\tau[Y],\E)=\E(Y)\) for every abelian \(\tau\)-sheaf \(\E\)
\cite[Tag 03AB]{stacks}, so that
\(H^n_\tau(Y,\E)=\operatorname{Ext}^n(\Z_\tau[Y],\E)\), the Ext group being
computed in \(\Ab\Sh(A_P^{\mathrm{op}},\tau)\) (apply
\(\Hom(\Z_\tau[Y],-)\) to an injective resolution of \(\E\)).
Fix a sign cover of \(R\) at \(x\).  For every \(X\), the universal
property \eqref{eq:universal} identifies \(h_{R_\pm}(X)\) with subsets of
\(h_R(X)\), with intersection \(h_{R_0}(X)\) by \eqref{eq:intersection}.
The elementary exact sequence
\(0\to\Z[U\cap U']\to\Z[U]\oplus\Z[U']\to\Z[U\cup U']\to0\) of free abelian
groups on two subsets of a set therefore gives an exact sequence of
presheaves
\[
 0\longrightarrow \Z[h_{R_0}]
  \longrightarrow \Z[h_{R_+}]\oplus\Z[h_{R_-}]
  \longrightarrow \Z[h_R]
\]
whose cokernel is \(\Z[h_R]/\Z[h_{R_+}\cup h_{R_-}]\).  This cokernel has
zero \(\tau\)-sheafification: a section of \(h_R\) over \(X\) is a map
\(f:R\to X\), and on the two members of the sign cover of \(X\) at
\(f^{\gp}(x)\), which are \(X\amalg_RR_\pm\) by \eqref{eq:basechange}, the
restrictions of \(f\) lie in \(h_{R_+}\) and \(h_{R_-}\) respectively.
Since sheafification is exact we obtain a short exact sequence of
\(\tau\)-sheaves
\[
 0\longrightarrow \Z_\tau[R_0]
  \longrightarrow \Z_\tau[R_+]\oplus\Z_\tau[R_-]
  \longrightarrow \Z_\tau[R]\longrightarrow0,
\]
and applying \(\operatorname{Ext}^\bullet(-,\E)\) the long exact sequence
\begin{equation}
\begin{split}
 0\to\E(R)\to\E(R_+)\oplus\E(R_-)\to\E(R_0)\to H^1_\tau(R,\E)\to\cdots
 \\
 \cdots\to H^{n-1}_\tau(R_0,\E)\to H^n_\tau(R,\E)
 \to H^n_\tau(R_+,\E)\oplus H^n_\tau(R_-,\E)\to H^n_\tau(R_0,\E)\to\cdots
\end{split}
\label{eq:mvlong}
\end{equation}
This is the Mayer--Vietoris sequence of \cite[Tag 0EVY]{stacks} for the
present situation.

\emph{Injectivity on the members of a sign cover.}  We prove the lemma by
induction on \(n\ge1\), simultaneously for all objects \(R\).  Consider the
statement \((\ast)_n\) that, for every sign cover of every object \(R\), the
map
\[
  H^n_\tau(R,\E)\longrightarrow
  H^n_\tau(R_+,\E)\oplus H^n_\tau(R_-,\E)
\]
is injective.  For \(n=1\), \((\ast)_1\) follows from \eqref{eq:mvlong} and the assumed
surjectivity of \eqref{eq:surj}, which makes the connecting map
\(\E(R_0)\to H^1_\tau(R,\E)\) zero.  For \(n>1\), if
\(H^{n-1}_\tau(-,\E)=0\) on all objects, then in particular
\(H^{n-1}_\tau(R_0,\E)=0\) and \eqref{eq:mvlong} gives \((\ast)_n\).

\review{This is probably the most important moment of the proof.  It appears correct, but is dismissed a little too informally.}

\emph{From the leaves to the root.}  Having established \((\ast)_n\)
(unconditionally for \(n=1\), and under the induction hypothesis for
\(n>1\)), let
\(\xi\in H^n_\tau(R,\E)\).  By locality of cohomology, \(\xi\) vanishes on
the members of some \(\tau\)-cover of \(R\).  By \eqref{eq:treesieve}, the
sieve generated by this family contains a tree family
\(\{R\to T_\ell\}\) (each \(T_\ell\) receives a map under \(R\) from a
member on which \(\xi\) vanishes), so \(\xi|_{T_\ell}=0\) for every leaf.  Now argue
upward through the tree: if \(M\) is an internal vertex with children
\(M_+,M_-\) forming a sign cover of \(M\), and \(\xi|_{M_+}=0=\xi|_{M_-}\),
then \(\xi|_M=0\) by \((\ast)_n\) applied to the sign cover of \(M\).
Induction on the height of a vertex gives \(\xi|_R=\xi=0\).  Hence
\(H^n_\tau(R,\E)=0\) for all \(R\), which completes the induction on \(n\).
\end{proof}

\begin{corollary}\label{cor:Gtau}
Let \(\G=\Lan_jF\) be the Kan extension of \S\ref{sec:kan}.  For every
object \(R\) of \(A_P\) and every \(n>0\), \(H^n_\tau(R,\G)=0\).
\end{corollary}

\begin{proof}
By Lemma~\ref{lem:mvsections}, \(\G\) is a \(\tau\)-sheaf whose maps
\eqref{eq:surj} are surjective.  Apply Lemma~\ref{lem:acyclic}.
\end{proof}

We now transfer this vanishing to the coarser topology \(\Jb\).  Some care
is needed, since the inclusion of \(\tau\)-sheaves into \(\Jb\)-sheaves is
not exact in general (cokernels involve different sheafifications); what
does transfer is injectivity of objects, and Corollary~\ref{cor:Gtau} makes
a \(\tau\)-injective resolution exact already as a complex of presheaves.

\begin{proposition}\label{prop:Gbir}
For \(\G=\Lan_jF\), every object \(R\) of \(A_P\) and every \(n>0\),
\(H^n_{\Jb}(R,\G)=0\).
\end{proposition}

\begin{proof}
Since \(\Jb\subset\tau\) (Corollary~\ref{cor:finer}), the identity of
\(A_P^{\mathrm{op}}\) is a morphism of sites
\((A_P^{\mathrm{op}},\tau)\to(A_P^{\mathrm{op}},\Jb)\) whose direct image is
the inclusion
\[
   \iota:\Ab\Sh(A_P^{\mathrm{op}},\tau)\longrightarrow\Ab\Sh(A_P^{\mathrm{op}},\Jb)
\]
and whose inverse image is \(\tau\)-sheafification, an exact functor
\cite[Tag 0EWK]{stacks}.  A right adjoint whose left adjoint is exact
preserves injectives \cite[Tag 015Z]{stacks}; hence \(\iota\) carries
injective \(\tau\)-sheaves to injective \(\Jb\)-sheaves.

Choose an injective resolution \(\G\to\I^\bullet\) in
\(\Ab\Sh(A_P^{\mathrm{op}},\tau)\).  By Corollary~\ref{cor:Gtau}, for every
object \(R\) the augmented complex of abelian groups
\begin{equation}
        0\longrightarrow\G(R)\longrightarrow\I^0(R)
         \longrightarrow\I^1(R)\longrightarrow\cdots\label{eq:complex}
\end{equation}
is exact.  So \(\G\to\I^\bullet\) is exact as a complex of presheaves, hence
exact as a complex of \(\Jb\)-sheaves (kernels of maps of sheaves are
computed in presheaves, and a presheaf image that is already a sheaf equals
the sheaf image).  Its terms \(\iota\I^k\) are injective \(\Jb\)-sheaves by
the preceding paragraph.  Thus \(\G\to\iota\I^\bullet\) is an injective
resolution of \(\G\) in \(\Ab\Sh(A_P^{\mathrm{op}},\Jb)\), and
\(H^n_{\Jb}(R,\G)\) is the \(n\)th cohomology of the same complex
\eqref{eq:complex}, which vanishes for \(n>0\).
\end{proof}

\section{Proof of Theorem~\ref{thm:main}}
\label{sec:unit}

Recall from \S\ref{sec:functors} that \(\pi^{-1}F\) is the \(\Jb\)-sheafification of the
left Kan extension \(\G=\Lan_jF\), and that \(\pi_*\) is restriction along
\(j:B_P\hookrightarrow A_P\).

\emph{\(\pi^{-1}F=\G\).}  By Lemma~\ref{lem:mvsections}, \(\G\) is already a
\(\Jb\)-sheaf, so no sheafification is needed and \(\pi^{-1}F=\G\).

\emph{The unit in degree zero.}  The unit \(F\to\pi_*\pi^{-1}F=\G\circ j\)
evaluated at \(L_a\) is the Kan extension unit \(F(L_a)\to\G(L_a)\), an
isomorphism by \eqref{eq:restrict}.

\emph{Vanishing of \(R^n\pi_*\) for \(n>0\).}  Let \(\G\to\I^\bullet\) be an
injective resolution in \(\Ab\Sh(A_P^{\mathrm{op}},\Jb)\).  By definition
\(\R\pi_*\G\) is represented by the complex of presheaves
\(\pi_*\I^\bullet=\I^\bullet\circ j\) on \(B_P^{\mathrm{op}}\).  Since
\(B_P^{\mathrm{op}}\) carries the trivial topology, evaluation at any object
\(L_a\) is an exact functor on \(\Ab\Sh(B_P^{\mathrm{op}})=\Ab\PSh(B_P^{\mathrm{op}})\),
so
\[
             (R^n\pi_*\G)(L_a)=H^n\bigl(\I^\bullet(L_a)\bigr)
             =H^n_{\Jb}(L_a,\G)=0
                                      \qquad(n>0)
\]
by Proposition~\ref{prop:Gbir}.  As the \(L_a\) are all the objects of
\(B_P\), the cohomology presheaves \(R^n\pi_*\G\) vanish for \(n>0\), and
\(R^0\pi_*\G=\pi_*\G\cong F\) by the previous paragraph.  Hence
\(F\to\R\pi_*\pi^{-1}F\) is a quasi-isomorphism.  This proves
Theorem~\ref{thm:main}, including the three ``more precise'' assertions.
\qed

\begin{remark}[Where the hypotheses are used]\label{rem:infinitesimal}
The infinitesimal lifting condition in the definition of a birhizomic cover
is never analysed: Proposition~\ref{prop:valuative} shows that \emph{any}
family of finitely presented maps through which every valuative map
factors is a \(\tau\)-cover, and this is all that is used.  Finite
presentation enters exactly once, in Step~3 of that proof, to factor
\(Q_i\to V_D\) through a finite stage of the colimit \eqref{eq:valcolim};
the colimit identity itself needs no finiteness.  Finite generation of
\(P\) is never used, which is why the compactness argument on \(\V(R)\) is
needed in place of a dimension count.  Finally, the argument of
Proposition~\ref{prop:Gbir} shows that \(\G\) is acyclic for every topology
coarser than \(\tau\), the birhizomic topology being one instance.
\end{remark}

\begin{remark}[Two small examples]
For \(R=\N^2\) and \(x=e_1-e_2\), neither sign lies in \(R\); the members
of the sign cover are \(\N^2\langle\pm(e_1-e_2)\rangle\), the two charts of
the blow-up of the origin, and \(R_0=\N^2\langle x,-x\rangle\simeq\N\) via
\((a,b)\mapsto a+b\).  None of these maps kills a nonzero element of
\(\N^2\), so \(\G\) takes the same value on all four monoids, as in the first
case of Lemma~\ref{lem:mvsections}.  For \(R=\N\) and \(x=1\) the sign
cover is \(\{\N\to\N,\ \N\to0\}\) with \(R_0=0\), and
\eqref{eq:mvsections} is \(0\to\G(\N)\to\G(\N)\oplus\G(0)\to\G(0)\to0\).
Deeper in a sign tree a leaf may kill nonzero elements of \(P\); this is
harmless, as Lemma~\ref{lem:zero} is applied at each internal vertex.
\end{remark}

\begin{thebibliography}{9}

\bibitem{CT}
B.~Conrad and M.~Temkin,
\emph{Non-Archimedean analytification of algebraic spaces},
J.~Algebraic Geom.\ 18 (2009), 731--788.

\bibitem{MW}
S.~Molcho and J.~Wise,
\emph{The logarithmic Picard group and its tropicalization},
Compos.\ Math.\ 158 (2022), 1477--1562.

\bibitem{niziol}
W.~Nizio{\l},
\emph{Toric singularities: log-blow-ups and global resolutions},
J.~Algebraic Geom.\ 15 (2006), 1--29.

\bibitem{ogus}
A.~Ogus,
\emph{Lectures on logarithmic algebraic geometry},
Cambridge Studies in Advanced Mathematics 178,
Cambridge University Press, 2018.

\bibitem{stacks}
The Stacks Project Authors,
\emph{The Stacks Project}, 2026.
Tags 00VC, 00VE, 00VF (Kan extension along a functor and its adjoint),
00ZC (the topology defined by a site),
03AB (free abelian sheaves),
015Z (exact left adjoints and injectives),
03P1 (locality of cohomology),
0EVY (Mayer--Vietoris on a site),
0EWK (comparing two topologies).

\end{thebibliography}

\end{document}
